\section*{Chapter \ref{ch:rc:reduced-form-credit} --- Reduced-Form Credit}

\begin{solution}{exo:rc:reduced-form-credit:1}
$\lambda \approx 0.0180/0.60 = 3.00\%$ a year; five-year survival $e^{-0.03\times5} = 86.07\%$,
a \omterm{def:rc:reduced-form-credit:pd}{probability of default} of 13.9\%.
\end{solution}

\begin{solution}{exo:rc:reduced-form-credit:2}
Each spread prices the average hazard up to its maturity. When spreads rise with maturity,
the later interval must carry a higher hazard than the average so far to lift the
average: forward hazards lie above average hazards, as forward rates lie above zero rates
on an upward-sloping curve.
\end{solution}

\begin{solution}{exo:rc:reduced-form-credit:3}
The seller pays $(1-R)N = \text{USD}~6$ million and loses its position, worth
$-\text{USD}~88\,133$ to it: the jump-to-default is $-6\,000\,000+88\,133 = -\text{USD}~5\,911\,867$,
the mirror image of the buyer's.
\end{solution}

\begin{solution}{exo:rc:reduced-form-credit:4}
A one-basis-point rise in the spread raises the protection buyer's value by about the
\omterm{def:rc:reduced-form-credit:risky}{risky annuity} times the notional times $10^{-4}$: $4.41\times10\,000\,000\times10^{-4} \approx
\text{USD}~4\,407$, close to the bucketed USD~4\,391. The difference is the fall of the
annuity itself as the hazard rises.
\end{solution}

\begin{solution}{exo:rc:reduced-form-credit:5}
At the 500 coupon the quote is below the coupon, so the buyer receives an upfront of
5.98\% of notional (an upfront of $-5.98\%$ paid by the buyer); at the 100 coupon the
buyer would pay 9.96\%.
\end{solution}

\begin{solution}{exo:rc:reduced-form-credit:6}
\omterm{def:rc:reduced-form-credit:cs01}{CS01} is a derivative: it describes small, continuous changes of spreads while the name
survives. Default is a jump of the default indicator: the protection pays a fixed loss and
the position's premium stream stops, whatever the spread was just before. Spreads blowing
out to thousands of basis points approach but never reach the default payoff, and a
linear extrapolation of \omterm{def:rc:reduced-form-credit:cs01}{CS01} overshoots or undershoots it depending on the position.
\end{solution}

\begin{solution}{exo:rc:reduced-form-credit:7}
\texttt{basis\_package()}: the bond-implied spread is 220 basis points, the basis
$-99.8$ basis points, and the jump-to-default with USD~100 million of protection
$100\text{m}\times(0.40-0.97)+100\text{m}\times0.60-88\,133\times10 = \text{USD}~2\,118\,667$.
\end{solution}

\begin{solution}{exo:rc:reduced-form-credit:8}
Two flaws. The \omterm{def:rc:reduced-form-credit:triangledef}{credit triangle} gives a hazard of $0.0120/0.60 = 2\%$ only for a recovery
of 40\%, and a flat curve; with this chapter's upward-sloping curve the one-year hazard is
1.00\%. And the probabilities in spreads are risk-neutral: they include the \omterm{def:rc:reduced-form-credit:premium}{default risk
premium}, so the market's real-world estimate of default is lower.
\end{solution}

\begin{solution}{pb:rc:reduced-form-credit:1}
\textbf{1.} 101.38 per 100.
\textbf{2.} 220 basis points.
\textbf{3.} $120-220 = -99.8$ basis points (to one decimal).
\textbf{4.} The bond must be funded and financed (repo haircuts, balance sheet), it may be
illiquid or in forced selling, the bond and swap may reference different deliverables,
and in 2008 funding costs dwarfed default risk.
\textbf{5.} The bond's yield over funding, less the protection coupon, plus the pull of the
basis to zero if it converges; the upfront paid on the protection is part of the cost.
\textbf{6.} USD~2\,118\,667.
\textbf{7.} USD~345\,107.
\textbf{8.} $-\text{USD}~2\,141$ (par for par) and $-\text{USD}~3\,453$ (market value) per basis
point.
\textbf{9.} 28.5\%.
\textbf{10.} Par for par hedges the face that the default swap delivers against, so the
package gains on default at any recovery; the market-value hedge is cheaper and leaves
less spread risk but loses on default below 28.5\% recovery. Funds that can tolerate
spread risk choose par for par.
\textbf{11.} The package marks down: the bond falls more than the protection gains,
through the negative \omterm{def:rc:reduced-form-credit:cs01}{CS01}, although the carry is unchanged if held to maturity.
\textbf{12.} The bond is financed in repo; a rise in haircuts or repo rates, or a withdrawn
line, forces a sale at the widest basis.
\textbf{13.} On default the buyer can deliver the cheapest deliverable obligation; settlement
by auction sets one final price for all, which reduces the option's value, but a package
long a bond trading above the auction price still earns the difference.
\textbf{14.} The seller may default together with the name or in a crisis: the protection is
worth least when it is needed. Clearing removes most of it for standard contracts.
\textbf{15.} Holders sell bonds to raise cash, the basis widens sharply, repo tightens, and the
package marks down while nothing has defaulted, as in the fourth quarter of 2008.
\textbf{16.} It is riskless on default and to maturity, but its value, funding and margin in
between are not, and a leveraged holder can be forced out before convergence.
\textbf{17.} Bank credit desks and relative-value funds with cheap, stable funding and room on
the balance sheet; the capital cost of holding the bond decides who can run it.
\textbf{18.} Funding and its term, the deliverables, the counterparty and clearing of the
protection, the recovery assumption, and the basis's history in stress.
\textbf{19.} Named result: \emph{the negative-basis package}: a basis of $-99.8$ basis points,
and with the par-for-par hedge a jump-to-default gain of USD~2\,118\,667 and a \omterm{def:rc:reduced-form-credit:cs01}{CS01} of
$-\text{USD}~2\,141$ per basis point on USD~100 million.
\textbf{20.} What it costs to fund and hold a bond compared with holding the same credit risk
unfunded, in a swap.
\end{solution}

\begin{solution}{iq:rc:reduced-form-credit:1}
Spread $\approx$ hazard $\times$ \omterm{def:rc:reduced-form-credit:pd}{loss given default}. At 200 basis points and 40\% recovery the
hazard is $0.02/0.6 = 3.33\%$ a year, a five-year survival of about $e^{-0.167} = 84.6\%$.
\iqlookfor{the formula, its assumptions (flat hazard, continuous premium) and a quick number.}
\end{solution}

\begin{solution}{iq:rc:reduced-form-credit:2}
Take par spreads for increasing maturities and a recovery; solve for a piecewise-constant
hazard interval by interval so that each swap reprices, on a given \omterm{def:rc:multi-curve-and-collateral-discounting:curves}{discount curve}. Things
that go wrong: inconsistent quotes needing negative hazards, stale quotes, the recovery
assumption, interpolation choices between \omterm{def:rc:curve-construction:pillar}{pillars}, and the sensitivity of forward hazards
to small quote changes.
\iqlookfor{the procedure, the positivity check and the role of recovery.}
\end{solution}

\begin{solution}{iq:rc:reduced-form-credit:3}
Fixed coupons of 100 and 500 basis points with an upfront, quarterly standard dates,
auction settlement of credit events and central clearing. Fixed coupons make contracts on
a name fungible so that offsetting trades net; the upfront must be computed identically by
both sides and by the clearing house, hence one model with agreed conventions.
\iqlookfor{fungibility, netting and why the conversion must be common.}
\end{solution}

\begin{solution}{iq:rc:reduced-form-credit:4}
Default is a jump, not a spread move: \omterm{def:rc:reduced-form-credit:cs01}{CS01} measures small moves while the name survives,
whereas jump-to-default is the loss if it defaults now. A book flat \omterm{def:rc:reduced-form-credit:cs01}{CS01} across names can be
very long or short default (a short-dated long protection against a long-dated short), and
concentration limits and capital rules measure jump-to-default separately.
\iqlookfor{the distinction and a CS01-flat, JTD-long example.}
\end{solution}

\begin{solution}{iq:rc:reduced-form-credit:5}
Spreads compensate for expected losses, for the uncertainty and correlation of defaults
(which cluster in bad times), for liquidity and for funding; risk-neutral probabilities
absorb all of it. Pricing and hedging use risk-neutral probabilities; provisioning and
capital use real-world estimates (through the cycle or point in time).
\iqlookfor{the \omterm{def:rc:reduced-form-credit:premium}{default risk premium} and the right measure per use.}
\end{solution}

\begin{solution}{iq:rc:reduced-form-credit:6}
Buy the bond and buy protection (a negative-basis package), sized par for par or on market
value; finance the bond in repo. It can go wrong through funding, a further widening of the
basis marked to market, deliverability, recovery and the protection seller's risk.
\iqlookfor{the package, the hedge ratio and the funding risk.}
\end{solution}
